% Generated draft: requires mathematical and source review.
\documentclass[11pt]{article}
\usepackage[margin=25mm]{geometry}
\usepackage{fontspec}
\setmainfont{Latin Modern Roman}
\usepackage{amsmath,amssymb,mathtools,unicode-math}
\setmathfont{Latin Modern Math}
\usepackage{enumitem,needspace,xurl,hyperref}
\setlength{\parindent}{0pt}
\setlength{\parskip}{5pt}
\setlength{\emergencystretch}{3em}
\newcommand{\N}{\mathbb{N}}
\newcommand{\Z}{\mathbb{Z}}
\newcommand{\Q}{\mathbb{Q}}
\newcommand{\R}{\mathbb{R}}
\newcommand{\C}{\mathbb{C}}
\newcommand{\F}{\mathbb{F}}
\newcommand{\A}{\mathbb{A}}
\newcommand{\PP}{\mathbb{P}}
\newcommand{\E}{\mathbb{E}}
\newcommand{\eps}{\varepsilon}
\newcommand{\dd}{\,\mathrm{d}}
\begin{document}
\section*{UnsolvedMath ID: 30005383}
\textbf{Source problem code:} catalogue entry
\subsection{Definitions and mathematical statement}
Write \(C(K)\) for the graded Khovanov chain complex denoted by the double-bracket notation in the supplied record, and write \([C(K)]_{\sim}\) for its class under the intended equivalence relation \(\sim\). The record describes an abelian group \(G\) of such classes, with multiplication induced by tensor product:
\[
[C(K)]_{\sim}\cdot[C(L)]_{\sim}
=[C(K)\otimes C(L)]_{\sim}.
\]
This formula records the intended operation; it does not supply the missing definition of the equivalence relation or establish closure, well-definedness, or the group axioms.

\textbf{Problem (conditional on the intended construction of \(G\)).} Determine the isomorphism type, as an abelian group, of \(G\). The supplied original question is: ``What can one say about the isomorphism type of \(G\)?''

\textbf{Unresolved construction details.} The supplied record does not specify the objects allowed as \(K\), the coefficient ring, the precise Khovanov complex variant, the grading and shift conventions, the equivalence relation, or the tensor-product base ring. In particular, the title's reference to concordance does not define the equivalence relation. These omissions prevent an unconditional definition of \(G\) from this record alone; they must be resolved from the relevant report passage before the problem has a fully precise formulation.

\textbf{Supplied structural context.} The dated curation describes \(G\) as established and asserts that it admits a surjection onto an infinite-rank free abelian group whose kernel is infinitely generated, while its full isomorphism type remains undetermined. This is retained as a supplied literature assessment, not as an independently verified theorem or a substitute for the missing construction.
\subsection{Short English statement}
Determine the abstract abelian-group structure of the group that the source associates with equivalence classes of graded Khovanov chain complexes under tensor product. The supplied record leaves the equivalence relation and several construction conventions unspecified, so this statement remains conditional on recovering those definitions. Its \texttt{partially\_solved} label and structural claims are supplied dated curation requiring independent review.
\subsection{Sources}
\begin{itemize}
\item Supplied local UnsolvedMath registry record: permanent ID \(30005383\), source problem code \texttt{catalogue entry}, title \emph{Isomorphism Type of the Khovanov-Complex Concordance Group}; collection: \emph{Oberwolfach Reports Open Problems}. The record gives proposed year \(2023\) and supplies no individual proposer attribution.
\item \emph{Morphisms in Low Dimensions} (2023), Oberwolfach Reports, pp.\ 215--259, DOI \texttt{10.4171/owr/2023/4}. Supplied source and literature-source label: Oberwolfach Report 4/2023. \url{https://doi.org/10.4171/owr/2023/4}. The record describes these reports as documenting mathematical workshops and their open-problem discussions; no precise page or complete passage defining \(G\) is included in the supplied record.
\item Schuetz, \emph{Integral Rasmussen invariant}, supplied literature reference: \url{https://arxiv.org/abs/2202.00445}. The record does not identify a particular result in this reference supporting the asserted structure of \(G\).
\item Literature assessment date: \texttt{2026-08-22}. The supplied dated curation labels the problem \texttt{partially\_solved} and states that \(G\) has a surjection onto an infinite-rank free abelian group with infinitely generated kernel, but that its full isomorphism type is undetermined. This assessment is supplied dated curation, not independently verified; the record expressly calls for independent verification of time-sensitive status and cited claims.
\item The supplied record labels the statement \texttt{corrected\_verified} and reports that it was checked against the official Oberwolfach report. This is a source-record verification claim, not verification performed for this draft.
\end{itemize}
\end{document}
